\section{THỰC TRẠNG HOẠT ĐỘNG QUẢN TRỊ TẠI FPT}

\subsection{Cơ cấu quản trị doanh nghiệp (Corporate Governance)}
Hệ thống quản trị của FPT được tổ chức theo mô hình công ty đại chúng niêm yết chuẩn mực gồm 3 định chế chính:
\begin{enumerate}
    \item \textbf{Đại hội đồng Cổ đông (ĐHĐCĐ):} Cơ quan quyết định cao nhất của công ty cổ phần, thông qua báo cáo tài chính, phân phối lợi nhuận, bầu HĐQT và Ban Kiểm soát.
    \item \textbf{Hội đồng Quản trị (HĐQT):} Cơ quan quản trị cao nhất giữa hai kỳ họp ĐHĐCĐ, chịu trách nhiệm định hướng chiến lược dài hạn và giám sát Ban Điều hành.
    \item \textbf{Ủy ban Kiểm toán / Ban Kiểm soát:} Giám sát tính hợp pháp và trung thực trong công tác kế toán, kiểm soát nội bộ và quản trị rủi ro.
\end{enumerate}

\subsection{Vai trò của Hội đồng Quản trị và tính độc lập trong giám sát}
\begin{itemize}
    \item \textbf{Chủ tịch HĐQT (Trương Gia Bình):} Định hình văn hóa, tầm nhìn chiến lược vĩ mô, xúc tiến các quan hệ đối tác quốc tế chiến lược (như với NVIDIA, các chính phủ).
    \item \textbf{Cơ cấu thành viên HĐQT:} Tỷ lệ thành viên độc lập và thành viên không điều hành theo quy định của Luật Doanh nghiệp 2020 và Thông tư số 116/2020/TT-BTC.
    \item \textbf{Các tiểu ban chuyên trách trực thuộc HĐQT:} Tiểu ban Lương thưởng, Tiểu ban Nhân sự và Bổ nhiệm, Tiểu ban Chiến lược.
\end{itemize}

\subsection{Phân định ranh giới giữa Quản trị (Governance) và Quản lý (Management)}
Bản chất sự khác biệt tại FPT:
\begin{table}[htbp]
\centering
\caption{So sánh vai trò Quản trị và Quản lý tại Tập đoàn FPT}
\label{tab:quan_tri_vs_quan_ly}
\begin{tabularx}{\textwidth}{p{3.2cm} X X}
\toprule
\textbf{Tiêu chí} & \textbf{Quản trị (Governance)} & \textbf{Quản lý (Management)} \\
 chủ thể & Hội đồng Quản trị (HĐQT) & Ban Điều hành (CEO, Ban Giám đốc) \\
Trọng tâm & Định hướng chiến lược, chính sách, giám sát & Thực thi kế hoạch kinh doanh, tổ chức vận hành \\
Khung thời gian & Dài hạn (3--10 năm) & Ngắn và trung hạn (hàng ngày, quý, năm) \\
Trách nhiệm giải trình & Chịu trách nhiệm trước ĐHĐCĐ & Chịu trách nhiệm trước HĐQT \\
\bottomrule
\end{tabularx}
\end{table}

\subsection{Đánh giá việc tuân thủ các chuẩn mực quản trị tốt}
Đối chiếu với \textit{Bộ nguyên tắc Quản trị Công ty theo thông lệ tốt nhất tại Việt Nam (CG Code)} do UBCKNN ban hành và chuẩn mực OECD:
\begin{itemize}
    \item \textbf{Điểm sáng:} Minh bạch thông tin, công bố báo cáo thường niên song ngữ Anh - Việt chất lượng cao, tiếp xúc nhà đầu tư chuyên nghiệp.
    \item \textbf{Điểm cần lưu ý:} Vẫn còn tỷ lệ giao dịch với các bên liên quan trong hệ sinh thái (như FPT Retail, Synnex FPT); tính độc lập thực sự của các thành viên độc lập trước ảnh hưởng quá lớn của Chủ tịch HĐQT.
\end{itemize}
